% Part: turing-machines
% Chapter: machines-computations
% Section: representing-tms

\documentclass[../../../include/open-logic-section]{subfiles}

\begin{document}

\olfileid{tur}{mac}{rep}
\olsection{ట్యూరింగ్ యంత్రాలకు ప్రాతినిధ్యం}

\begin{explain}
ట్యూరింగ్ యంత్రాలను \emph{స్థితి రేఖాచిత్రాల}తో దృశ్యరూపంలో చూపవచ్చు.
బాణాలతో కలిపిన స్థితి గడులతో ఈ రేఖాచిత్రాలు ఏర్పడతాయి. సహజంగానే, ప్రతి
స్థితి గడి యంత్రపు ఒక స్థితికి ప్రాతినిధ్యం వహిస్తుంది. ప్రతి బాణం ఆ స్థితి
నుండి అమలు చేయగల ఒక నిర్దేశానికి ప్రాతినిధ్యం వహిస్తుంది; ఆ నిర్దేశపు
వివరాలను సంబంధిత బాణం పైనా కిందా వ్రాస్తారు. $q_0$, $q_1$ అనే రెండు అంతర్గత
స్థితులు, ఒకే నిర్దేశం ఉన్న కింది యంత్రాన్ని పరిశీలించండి:
\[
\begin{tikzpicture}[->,>=stealth',shorten >=1pt,auto,node distance=2.8cm,
                    semithick]
  \tikzstyle{every state}=[fill=none,draw=black,text=black]

  \node[initial,state]   (A)              {$q_0$};
  \node[state]   (B) [right of=A] {$q_1$};

  \path (A) edge  node {\TMtrans{\TMblank}{\TMstroke}{\TMright}} (B);
\end{tikzpicture}
\]
ట్యూరింగ్ యంత్రానికి చదువు--వ్రాత శీర్షమూ, దానిపై ఇన్‌పుట్ వ్రాసి ఉన్న
టేపూ ఉంటాయని గుర్తుచేసుకోండి. ఈ నిర్దేశాన్ని ఇలా చదవవచ్చు: \emph{స్థితి
$q_0$లో ఒక~$\TMblank$ను చదువుతుంటే, ఒక~$\TMstroke$ను వ్రాసి, కుడికి కదిలి,
స్థితి $q_1$కు వెళ్లాలి}. ఇది $\tuple{q_0, \TMblank}$ను
$\tuple{q_1, \TMstroke, \TMright}$కు పంపే సంక్రమణ ప్రమేయానికి సమానం.
\end{explain}

\begin{ex}
\emph{సరి యంత్రం}: టేపుపై సరి సంఖ్యలో $\TMstroke$లు ఉంటే, అప్పుడే మాత్రమే
కింది ట్యూరింగ్ యంత్రం ఆగుతుంది (టేపుపై మొదటి $\TMblank$కు ముందు అన్ని
$\TMstroke$లు వస్తాయని ఊహించాం).
\[
\begin{tikzpicture}[->,>=stealth',shorten >=1pt,auto,node distance=2.8cm,
                    semithick]
  \tikzstyle{every state}=[fill=none,draw=black,text=black]

  \node[initial,state]         (A)              {$q_0$};
  \node[state]         (B) [right of=A] {$q_1$};

  \path (A) edge [bend left] node {\TMtrans{\TMstroke}{\TMstroke}{\TMright}} (B)
        (B) edge [loop above] node {\TMtrans{\TMblank}{\TMblank}{\TMright}} (B)
            edge [bend left] node {\TMtrans{\TMstroke}{\TMstroke}{\TMright}} (A);
\end{tikzpicture}
\]

ఈ స్థితి రేఖాచిత్రం కింది సంక్రమణ ప్రమేయానికి అనుగుణంగా ఉంటుంది:
\begin{align*}
\delta(q_0, \TMstroke) & = \tuple{q_1, \TMstroke, \TMright},\\
\delta(q_1, \TMstroke) & = \tuple{q_0, \TMstroke, \TMright},\\
\delta(q_1, \TMblank)  & = \tuple{q_1, \TMblank, \TMright}
\end{align*}
\end{ex}

\begin{explain}
ఇన్‌పుట్‌లో సరి సంఖ్యలో గీతలు ఉన్నప్పుడే పై యంత్రం ఆగుతుంది. లేకపోతే యంత్రం
(సిద్ధాంతపరంగా) అనంతకాలం పనిచేస్తూనే ఉంటుంది. ఏ యంత్రం, ఇన్‌పుట్ ఇచ్చినా,
అవుట్‌పుట్‌ను నిర్ణయించడానికి ఆ యంత్రపు \emph{స్థితివిన్యాసాలను} వరుసగా
అనుసరించవచ్చు. స్థితివిన్యాసాలకు ఆచారబద్ధ నిర్వచనాన్ని తరువాత ఇస్తాం.
ప్రస్తుతానికి, యంత్రం పనిచేసే సమయంలో ఏ క్షణంలోనైనా దాని స్థితిని చూపే
రేఖాచిత్రాల వరుసగా స్థితివిన్యాసాలను అంతర్బోధాత్మకంగా భావించవచ్చు.
స్థితివిన్యాసాలు టేపు విషయం, యంత్రపు స్థితి, చదువు--వ్రాత శీర్షం ఉన్న స్థానాన్ని
చూపుతాయి.

సరి యంత్రాన్ని నాలుగు $\TMstroke$ల ఇన్‌పుట్‌తో ప్రారంభించినప్పుడు దాని
స్థితివిన్యాసాలను అనుసరిద్దాం. ఈ సందర్భంలో యంత్రం ఆగుతుందని ఆశిస్తాం. తరువాత
మూడు $\TMstroke$ల ఇన్‌పుట్‌పై యంత్రాన్ని నడుపుతాం; అప్పుడు అది ఎప్పటికీ
నడుస్తుంది.

యంత్రం స్థితి~$q_0$లో, అత్యంత ఎడమ $\TMstroke$ను చదువుతూ ప్రారంభమవుతుంది.
యంత్రపు ప్రారంభ స్థితివిన్యాసాన్ని ఇలా చూపవచ్చు:
\[
\TMendtape \TMstroke_0 \TMstroke \TMstroke \TMstroke \TMblank \ldots
\]
పై స్థితివిన్యాసం సూటిగా ఉంది. కనిపిస్తున్నట్లుగా, యంత్రం స్థితి~$q_0$లో
అత్యంత ఎడమ $\TMstroke$ను చదువుతూ ప్రారంభమవుతుంది. మొదటి $\TMstroke$పై ఉన్న
స్థితి పేరు యొక్క ఉపలిపి దీనిని సూచిస్తుంది. ఈ బిందువులో వర్తించే నిర్దేశం
$\delta(q_0, \TMstroke) = \tuple{q_1, \TMstroke, \TMright}$; కాబట్టి యంత్రం
టేపుపై కుడికి కదిలి స్థితి~$q_1$కు మారుతుంది.
\sourcecorrection{OLTETURMAC-001}{ఆంగ్ల మూలం ఈ స్థితివిన్యాసంలో యంత్రం ``state one''లో మొదలవుతుందని చెప్పింది; కానీ ముందరి వాక్యం, గీతపైని ఉపలిపి, వర్తించే నిర్దేశం అన్నీ ప్రారంభ స్థితి $q_0$ అని చూపుతున్నాయి. అందుకు అనుగుణంగా $q_0$గా స్పష్టం చేశాం.}
\[
\TMendtape \TMstroke \TMstroke_1 \TMstroke \TMstroke \TMblank \ldots
\]
యంత్రం ఇప్పుడు స్థితి~$q_1$లో ఒక~$\TMstroke$ను చదువుతోంది కనుక,
$\delta(q_1, \TMstroke) = \tuple{q_0, \TMstroke, \TMright}$ అనే నిర్దేశాన్ని
``అనుసరించాలి.'' దీని ఫలితంగా వచ్చే స్థితివిన్యాసం
\[
\TMendtape \TMstroke \TMstroke \TMstroke_0 \TMstroke \TMblank \ldots
\]
యంత్రం కొనసాగేకొద్దీ, నియమాలు మళ్లీ అదే క్రమంలో వర్తించి కింది రెండు
స్థితివిన్యాసాలు వస్తాయి:
\[
\TMendtape \TMstroke \TMstroke \TMstroke \TMstroke_1 \TMblank \ldots
\]
\[
\TMendtape \TMstroke \TMstroke \TMstroke \TMstroke \TMblank_0 \ldots
\]
యంత్రం ఇప్పుడు స్థితి~$q_0$లో ఒక~$\TMblank$ను చదువుతోంది. స్థితి రేఖాచిత్రాన్ని
బట్టి, అమలు చేయడానికి నిర్దేశమేదీ లేదని సులభంగా చూడవచ్చు; అందువల్ల యంత్రం
ఆగిపోయింది. అంటే ఇన్‌పుట్ అంగీకరించబడింది.

తరువాత యంత్రాన్ని మూడు $\TMstroke$ల ఇన్‌పుట్‌తో ప్రారంభిద్దాం. అవే నిర్దేశాలు
అమలవుతాయి కనుక మొదటి కొన్ని స్థితివిన్యాసాలు పోలి ఉంటాయి; టేపు ఇన్‌పుట్‌లో
మాత్రమే చిన్న తేడా ఉంటుంది:
\[
\TMendtape \TMstroke_0 \TMstroke \TMstroke \TMblank \ldots
\]
\[
\TMendtape \TMstroke \TMstroke_1 \TMstroke \TMblank \ldots
\]
\[
\TMendtape \TMstroke \TMstroke \TMstroke_0 \TMblank \ldots
\]
\[
\TMendtape \TMstroke \TMstroke \TMstroke \TMblank_1 \ldots
\]
యంత్రం ఇప్పుడు అన్ని $\TMstroke$లను దాటి వెళ్లి, స్థితి~$q_1$లో ఒక~$\TMblank$ను
చదువుతోంది. రేఖాచిత్రంలో చూపినట్లుగా $\delta(q_1, \TMblank) =
\tuple{q_1, \TMblank, \TMright}$ రూపంలోని నిర్దేశం ఉంది. టేపు కుడివైపు
అనంతంగా $\TMblank$లతో నిండి ఉంటుంది కనుక, యంత్రం స్థితి~$q_1$లోనే ఉండి ఇంకా
ఇంకా కుడికి కదులుతూ ఈ నిర్దేశాన్ని \emph{ఎప్పటికీ} అమలు చేస్తూనే ఉంటుంది.
యంత్రం ఎన్నటికీ ఆగదు; ఇన్‌పుట్‌ను అంగీకరించదు.
\end{explain}

\begin{explain}
అన్ని యంత్రాలూ ఆగవని గమనించడం ముఖ్యం. ఆగడం అంటే యంత్రానికి అమలు చేయడానికి
నిర్దేశాలు అయిపోవడమే అయితే, ప్రతి స్థితిలో ప్రతి సంకేతానికి ఒక నిష్క్రమణ బాణం
ఉండేలా చేయడం ద్వారా ఎన్నటికీ ఆగని యంత్రాన్ని నిర్మించవచ్చు. $q_0$లో ఒక
$\TMblank$ను చదివే సందర్భానికి ఒక నిర్దేశాన్ని చేర్చి, సరి యంత్రం అనంతకాలం
నడిచేలా మార్చవచ్చు.
\end{explain}

\begin{ex}
\[
\begin{tikzpicture}[->,>=stealth',shorten >=1pt,auto,node distance=2.8cm,
                    semithick]
  \tikzstyle{every state}=[fill=none,draw=black,text=black]

  \node[initial,state] (A)              {$q_0$};
  \node[state]         (B) [right of=A] {$q_1$};

  \path
  (A) edge [bend left] node {\TMtrans{\TMstroke}{\TMstroke}{\TMright}} (B)
      edge [loop above] node {\TMtrans{\TMblank}{\TMblank}{\TMright}} (A)
  (B) edge [loop above] node {\TMtrans{\TMblank}{\TMblank}{\TMright}} (B)
      edge [bend left] node {\TMtrans{\TMstroke}{\TMstroke}{\TMright}} (A);
\end{tikzpicture}
\]
\end{ex}

\begin{explain}
ట్యూరింగ్ యంత్రాలకు ప్రాతినిధ్యం వహించడానికి యంత్ర పట్టికలు మరొక పద్ధతి.
యంత్ర పట్టికల్లో టేపు వర్ణమాలను $x$-అక్షంపైనా, యంత్ర స్థితుల సమితిని
$y$-అక్షంపైనా చూపుతారు. పట్టికలో ప్రతి స్థితి, సంకేతం కలిసే గడిలో నిర్దేశపు
మిగిలిన భాగాన్ని---కొత్త స్థితి, కొత్త సంకేతం, కదలిక దిశను---వ్రాస్తారు.
యంత్రం ఏ స్థితిలో, ఏ సంకేతానికి ఆగుతుందో నిర్ణయించడం యంత్ర పట్టికలతో సులభం.
పట్టికలో ఖాళీ ఉన్న ప్రతి చోటూ యంత్రం ఆగగల ఒక బిందువు. యంత్రం ఆగే బిందువులు
ఎల్లప్పుడూ వెంటనే స్పష్టంగా కనిపించని స్థితి రేఖాచిత్రాలు, నిర్దేశ సమితులకు
భిన్నంగా, యంత్ర పట్టికలోని ఖాళీలను కనుగొని ఏ ఆగే బిందువునైనా వెంటనే
గుర్తించవచ్చు.
\end{explain}

\begin{ex}
సరి యంత్రానికి యంత్ర పట్టిక:
\[
\centering
\begin{tabular}{lllll}
\cline{2-4}
\multicolumn{1}{l|}{}      & \multicolumn{1}{c|}{$\TMblank$}
& \multicolumn{1}{c|}{$\TMstroke$}     & \multicolumn{1}{c|}{$\TMendtape$}   &  \\ \cline{1-4}
\multicolumn{1}{|l|}{$q_0$} & \multicolumn{1}{l|}{}
& \multicolumn{1}{l|}{\TMtrans{\TMstroke}{q_1}{\TMright}} &
\multicolumn{1}{l|}{\phantom{\TMtrans{\TMstroke}{q_1}{\TMright}}} &  \\ \cline{1-4}
\multicolumn{1}{|l|}{$q_1$} & \multicolumn{1}{l|}{\TMtrans{\TMblank}{q_1}{\TMright}}
& \multicolumn{1}{l|}{\TMtrans{\TMstroke}{q_0}{\TMright}} &
\multicolumn{1}{l|}{\phantom{\TMtrans{\TMstroke}{q_1}{\TMright}}} &  \\ \cline{1-4}
\end{tabular}
\]
మనకు కనిపిస్తున్నట్లుగా, యంత్రం స్థితి~$q_0$లో ఒక~$\TMblank$ను చదివేటప్పుడు
ఆగుతుంది.
\end{ex}

\begin{explain}
ఇంతవరకు ఇన్‌పుట్‌ను చదివి అంగీకరించే యంత్రాలను మాత్రమే పరిశీలించాం. అయితే,
ట్యూరింగ్ యంత్రాలు చదవగలవు, వ్రాయగలవు కూడా. అలాంటి యంత్రానికి (ఇలాంటివి
చాలాచాలా ఉన్నప్పటికీ) ఒక ఉదాహరణ \emph{ద్విగుణక యంత్రం}. టేపుపై $n$
$\TMstroke$ల దిమ్మతో ప్రారంభించినప్పుడు, ద్విగుణక యంత్రం $2n$ $\TMstroke$ల
దిమ్మను అవుట్‌పుట్‌గా ఇస్తుంది.
\end{explain}

\begin{ex}
  \ollabel{ex:doubler} ద్విగుణక యంత్రాన్ని నిర్మించే ముందు సమస్యను పరిష్కరించే
  ఒక \emph{వ్యూహాన్ని} రూపొందించడం ముఖ్యం. మనం రూపొందించిన యంత్రం తాను ఎన్ని
  $\TMstroke$లను చదివిందో జ్ఞాపకం ఉంచుకోలేదు కనుక, టేపుపై ఉన్న అన్ని
  $\TMstroke$లను గుర్తుపెట్టుకునే పద్ధతి కావాలి. ఒక~$\TMblank$తో అవుట్‌పుట్‌ను
  ఇన్‌పుట్ నుండి వేరు చేయడం అలాంటి ఒక పద్ధతి. అప్పుడు యంత్రం ఇన్‌పుట్‌లోని
  మొదటి $\TMstroke$ను చెరిపి, మిగిలిన ఇన్‌పుట్ మీదుగా వెళ్లి, ఒక $\TMblank$ను
  వదిలి, రెండు కొత్త $\TMstroke$లను వ్రాయగలదు. తరువాత యంత్రం వెనక్కి వెళ్లి
  ఇన్‌పుట్‌లోని రెండవ $\TMstroke$ను కనుగొని దానిని కూడా ద్విగుణీకరిస్తుంది.
  ఇన్‌పుట్‌లోని ప్రతి ఒక్క $\TMstroke$కు అవుట్‌పుట్‌లో రెండు $\TMstroke$లను
  వ్రాస్తుంది. యంత్రం ముందుకు సాగేకొద్దీ ఇన్‌పుట్‌ను చెరిపివేయడం వల్ల ఏ
  $\TMstroke$నూ విడిచిపెట్టలేదనీ, దేనినీ రెండుసార్లు ద్విగుణీకరించలేదనీ
  నిర్ధారించవచ్చు. ఇన్‌పుట్ మొత్తం చెరిగిపోయినప్పుడు టేపుపై $2n$
  $\TMstroke$లు మిగులుతాయి. ఫలితంగా వచ్చే ట్యూరింగ్ యంత్రపు స్థితి
  రేఖాచిత్రాన్ని \olref{fig:doubler}లో చూపాం.
  \begin{figure}
\[
\begin{tikzpicture}[->,>=stealth',shorten >=1pt,auto,node distance=2.8cm,
                    semithick]
  \tikzstyle{every state}=[fill=none,draw=black,text=black]

  \node[initial,state] (1)              {$q_0$};
  \node[state]         (2) [right of=1] {$q_1$};
  \node[state]         (3) [right of=2] {$q_2$};
  \node[state]         (4) [below of=3] {$q_3$};
  \node[state]         (5) [left of=4]  {$q_4$};
  \node[state]         (6) [left of=5]  {$q_5$};

  \path (1) edge node {\TMtrans{\TMstroke}{\TMblank}{\TMright}} (2)
    (2) edge [loop above] node {\TMtrans{\TMstroke}{\TMstroke}{\TMright}} (2)
      edge node {\TMtrans{\TMblank}{\TMblank}{\TMright}} (3)
    (3) edge [loop above] node {\TMtrans{\TMstroke}{\TMstroke}{\TMright}} (3)
        edge node {\TMtrans{\TMblank}{\TMstroke}{\TMright}} (4)
    (4) edge [loop below] node {\TMtrans{\TMblank}{\TMstroke}{\TMleft}} (4)
        edge node {\TMtrans{\TMstroke}{\TMstroke}{\TMleft}} (5)
    (5) edge [loop below]  node {\TMtrans{\TMstroke}{\TMstroke}{\TMleft}} (5)
        edge              node {\TMtrans{\TMblank}{\TMblank}{\TMleft}} (6)
    (6) edge [loop below] node {\TMtrans{\TMstroke}{\TMstroke}{\TMleft}} (6)
        edge              node {\TMtrans{\TMblank}{\TMblank}{\TMright}} (1);
\end{tikzpicture}
\]
\caption{ఒక ద్విగుణక యంత్రం}
\ollabel{fig:doubler}
\end{figure}
\end{ex}

\begin{prob}
ఏదైనా ఒక ఇన్‌పుట్‌ను ఎంచుకుని, \olref[tur][mac][rep]{ex:doubler}లోని ద్విగుణక
యంత్రపు స్థితివిన్యాసాలను వరుసగా అనుసరించండి.
\end{prob}

\begin{prob}
$\{\TMendtape,\TMblank, A, B\}$ అనే వర్ణమాల గల ఒక ట్యూరింగ్ యంత్రాన్ని
రూపొందించండి. $A$ల సంఖ్య $B$ల సంఖ్యకు సమానమై, \emph{అలాగే} అన్ని $A$లు
అన్ని $B$లకు ముందు వచ్చే ఏ $A$, $B$ల సంకేతమాలనైనా అది అంగీకరించాలి, అంటే
దానిపై ఆగాలి. $A$ల సంఖ్య $B$ల సంఖ్యకు సమానం కాని లేదా అన్ని $A$లు అన్ని
$B$లకు ముందు రాని ఏ సంకేతమాలనైనా అది తిరస్కరించాలి, అంటే దానిపై ఆగకూడదు.
(ఉదాహరణకు, యంత్రం $AABB$, $AAABBB$లను అంగీకరించాలి; కానీ $AAB$,
$AABBAABB$ రెండింటినీ తిరస్కరించాలి.)
\end{prob}

\begin{prob}
$\{\TMendtape,\TMblank, A, B\}$ అనే వర్ణమాల గల ఒక ట్యూరింగ్ యంత్రాన్ని
రూపొందించండి. $A$, $B$లతో కూడిన ఏ సంకేతమాల~$\alpha$నైనా ఇన్‌పుట్‌గా తీసుకుని,
దానిని నకలు చేసి $\alpha\alpha$ రూపంలోని అవుట్‌పుట్‌ను ఉత్పత్తి చేయాలి.
(ఉదాహరణకు, ఇన్‌పుట్ $ABBA$ అయితే అవుట్‌పుట్ $ABBAABBA$ కావాలి.)
\end{prob}

\begin{prob}
\emph{వర్ణక్రమంలో ఉందా?}: $\{\TMendtape,\TMblank, A, B\}$ అనే వర్ణమాల గల
ఒక ట్యూరింగ్ యంత్రాన్ని రూపొందించండి. $A$, $B$ల పరిమిత క్రమాన్ని ఇన్‌పుట్‌గా
ఇచ్చినప్పుడు, అన్ని $A$లు అన్ని $B$లకు ఎడమవైపు ఉన్నాయో లేదో అది తనిఖీ చేయాలి.
ఇన్‌పుట్ సంకేతమాలను టేపుపై అలాగే ఉంచి, సంకేతమాల ``వర్ణక్రమంలో'' ఉంటే ఆగాలి;
లేకపోతే ఎప్పటికీ చక్రంలో తిరగాలి.
\end{prob}

\begin{prob}
\emph{వర్ణక్రమీకరణ యంత్రం}: $\{\TMendtape,\TMblank, A, B\}$ అనే వర్ణమాల గల
ఒక ట్యూరింగ్ యంత్రాన్ని రూపొందించండి. $A$, $B$ల పరిమిత క్రమాన్ని ఇన్‌పుట్‌గా
తీసుకుని, అన్ని $A$లు అన్ని $B$లకు ఎడమవైపు ఉండేలా వాటిని పునర్వ్యవస్థీకరించాలి.
(ఉదాహరణకు, $BABAA$ అనే క్రమం $AAABB$ కావాలి; $ABBABB$ అనే క్రమం $AABBBB$
కావాలి.)
\end{prob}

\end{document}
